\section{Representation Theory}\label{sec:rep-theory}

\subsection{Group Modules and Irreducible Representations}

The aim of this section if to serve as a review of the basics of representation theory.
This will not be an exhaustive review by any means.
For a more thorough mathematical treatment of the representation of finite groups we direct the reader to \cite{sagan2001symmetric}.

Let $G$ denote a finite group. We will denote the identity element of $G$ by $e$ to avoid confusion. 
A $G$-module is a finite dimensional vector space $V$ equipped with a group action $\cdot : G \times V \to V$ satisfying 
\begin{enumerate}
    \item \textbf{(Identity)} For all vectors $\ket{v} \in V$, the equality $e \cdot \ket{v} = \ket{v}$ is satisfied.
    \item \textbf{(Linearity)} For all linear combinations $\alpha \ket{v} + \beta \ket{w}$ and all group elements $g \in G$, the equality $g \cdot (\alpha \ket{v}+\beta\ket{b}) = \alpha (g\cdot \ket{v}) + \beta (g \cdot \ket{w})$ is satisfied.
    \item \textbf{(Compatability)} For all group elements $g,h \in G$ and vectors $\ket{v} \in V$, the equality $(gh) \cdot \ket{v} = g\cdot (h \cdot \ket{v})$ is satisfied.
\end{enumerate}
The linearity axiom implies the existence of a map $\pi: G \to \mathrm{End}(V)$ mapping the group element $g$ to the linear endomorphism $\ket{v} \mapsto g\cdot \ket{v}$. 
The identity and compatibility axioms imply that $\pi$ is actually a group homomorphism. 

\begin{lemma}
    We call a group homomorphism $\pi : G \to \mathrm{End}(V)$ a representation of $G$. If $\pi: G \to \mathrm{End}(V)$ is a representation, then $V$ is a $G$-module under 
    \begin{align}
        g \cdot \ket{v} = \pi(G)\ket{v}.
    \end{align}
    Moreover, if $V$ is a $G$-module, then the map $\pi(g)(\ket{v}) = g\cdot \ket{v}$ is a representation of $G$.
\end{lemma}
In light of this result, it is somewhat standard to conflate representation of $G$ and $G$-modules. 
As with all linear maps, representations are typically described by how each element of $G$ acts on each basis vector of $V$.
% Alternatively, one can describe a representation by first describing how the group $G$ acts on some set $X$ and then extending the action to the free vector space $\mathbb{C}[X]$ by linearity.

Sometimes, a representation can be reduced to a direct sum of smaller representations in the sense that 
\begin{align}
    \pi(g) = \pi_0(g) \oplus \pi_1(g) \oplus \dots \oplus \pi_k(g)
\end{align}
for each $\pi_i$ a representation.
In this case, we write $\pi = \pi_0 \oplus \dots \oplus \pi_k$.
This corresponds to a decomposition of $G$-modules $V = V_0 \oplus \dots \oplus V_k$ where each $V_i$ is a $G$-module.
If no such decomposition exists, then we call the representation (or $G$-module) irreducible.
Irreducible representations are typically referred to by the abbreviation ``irrep''.

\begin{theorem}[Maschke's Theorem \cite{sagan2001symmetric} Theorem 1.5.3] \label{thm:maschkes-theorem}
    Every representation $\pi : G \to \mathrm{End}(V)$ decomposes into a direct sum
    \begin{align}
        \pi \cong m_0 \pi_0 \oplus m_1 \pi_1 \oplus \dots \oplus m_k \pi_k
    \end{align}
    where $\pi_i$ and $\pi_j$ are distinct irreducible representations of $G$. Here $m_i \pi_i \cong \bigoplus_{j=0}^{m_i - 1} \pi_i$.

    Similarly, if $V$ is a finite dimensional $G$-module, then 
    \begin{align}
        V = m_0 V_0 \oplus \dots \oplus m_k V_k
    \end{align}
    where $V_i$ and $V_j$ are distinct irreducible $G$-modules.
\end{theorem}
Maschke's theorem requires two important caveats. 
Firstly, we must be working over an algebraically closed field. 
Secondly, it must be possible to average over the group, so an analog of Maschke's theorem holds for compact Lie groups.

\subsubsection{Schur's Lemma}

Suppose that $V$ is a $G$-module. 
We let $\mathrm{End}_G(V)$ denote all of the elements $T \in \mathrm{End}_G(V)$ such that for all $g \in G$
\begin{align}
    T \pi(g) = \pi(g)T.
\end{align}
Elements of $T$ are said to be invariant under the action of $G$ or $G$-linear.
This condition ends up being rather stringent.

\begin{lemma}[Schur's Lemma \cite{sagan2001symmetric} Theorem 1.6.5]\label{lem:schurs-lemma}
    Let $V$ and $W$ be $G$-modules and let $\pi_V$ and $\pi_W$ be the associated representations.
    We define $\mathrm{Hom}_G(V, W)$ by $T \in \mathrm{Hom}_G(V, W)$ if for all $g \in G$
    \begin{align}
        T \pi_V(g) = \pi_W(g) T.
    \end{align}
    If $V$ and $W$ are irreducible, then exactly one of the following holds.
    \begin{enumerate}
        \item The $G$-modules $V$ and $W$ are not isomorphic and $\mathrm{Hom}_G(V,W) = 0$.
        \item The $G$-modules $V$ and $W$ are isomorphic and $\mathrm{Hom}_G(V,W) \cong \mathbb{C}$ in the sense that $\mathrm{End}_G(V) = \{ \lambda I_V \mid \lambda \in \mathbb{C} \}$.
    \end{enumerate}
\end{lemma}

\begin{cor}[\cite{sagan2001symmetric} Theorem 1.7.9]
    Let $V$ be a $G$-module with such that
    \begin{align}
        V = m_0 V_0 \oplus \dots \oplus m_k V_k
    \end{align}
    with $V_i$ and $V_j$ distinct irreducible modules.
    Let $d_i$ be the dimension of $V_i$.
    Then the following hold.
    \begin{enumerate}
        \item The dimension of $V$ is given by $\mathrm{dim}(V) = m_0 d_0 + \dots + m_k d_k$.
        \item \label{item:artin-wedderburn} The endomorphism algebra of $V$ is isomorphic to the direct sum
        \begin{align}
            \mathrm{End}_G(V) \cong \bigoplus_{i = 0}^k \mathrm{End}(\mathbb{C}^{m_i}).
        \end{align}
        \item The endomorphism algebra $\mathrm{End}_G(V)$ has dimension $m_0^2 + m_1^2 + \dots + m_k^2$.
    \end{enumerate}
\end{cor}

\begin{proof}
    We will only include a proof of \Cref{item:artin-wedderburn} as it is the most important.
    From the decomposition of $V$, we get the equality
    \begin{align}
        \mathrm{End}_G(V) \cong \bigoplus_{i,j = 0}^k \mathrm{Hom}_G(m_iV_i,m_jV_j).
    \end{align}
    The elements of $\mathrm{Hom}_G(m_i V_i, m_j V_j)$ can be thought of as $m_i \times m_j$ matrices whose entries are in $\mathrm{Hom}_G(V_i, V_j)$. 
    However, if $i \neq j$, the by Schur's lemma we know that $\mathrm{Hom}_G(V_i, V_j)$ only contains the zero map. 
    Thus, all of the terms in the direct sum for $i \neq j$ are zero and can be ignored. 
    When $i = j$, the elements of $\mathrm{Hom}_G(V_i,V_i) = \mathrm{End}_G(V_i)$ are all multiples of the identity map. 
    So, the elements $g_i \in \mathrm{End}_G(m_iV_i)$ can be written in the form
    \begin{align}
        g_i = M_i \otimes I_{V_i},
    \end{align}
    where $M_i \in \mathrm{End}(\mathbb{C}^{m_i})$ is an $m_i \times m_i$ matrix with complex entries. 
    % Moreover, any such matrix $M_i$ defines an element $g_i = M_i \otimes I_{V_i} \in \mathrm{End}_G(m_i V_i)$.
    A routine argument shows $\mathrm{End}_G(m_iV_i)$ is isomorphic to the matrix algebra $\mathrm{End}(\mathbb{C}^{m_i})$ under the map $g_i \mapsto M_i$.
    Finally, we conclude that $\mathrm{End}_G(V) \cong \bigoplus_{i = 0}^k \mathrm{End}_G(m_i V_i)$ and hence
    \begin{align}
        \mathrm{End}_G(V) \cong \bigoplus_{i=0}^k \mathrm{End}(\mathbb{C}^{m_i}).
    \end{align}
\end{proof}

Concretely, Schur's lemma and its corollary impose a relatively stringent structure on the elements of $\mathrm{End}_G(V)$. 
For our purposes the most important property is that if $T \in \mathrm{End}_G(V)$ and $\ket{v} \in m_i V_i$ is supported entirely in copies of the $V_i$ module, then $T\ket{v} \in m_i V_i$ is also supported entirely in copies of $V_i$. 
The following result guarantees that the twirled error channel that we aim to study is invariant under the action of the Clifford group.

\begin{lemma}[Twirling] \label{lem:twirling}
    Suppose $V$ is a finite dimensional $G$-module with associated representation for $G$ finite. Given any $T \in \mathrm{End}(V)$ (notably $T$ need not be $G$-linear), the twirled map
    \begin{align}
        T_G = \frac{1}{|G|} \sum_{g \in G} \pi(g)^{-1} T \pi(g)
    \end{align}
    is an element of $\mathrm{End}_G(V)$.
\end{lemma}

\begin{proof}
    Fix $g \in G$. Conjugating $T_G$ by $\pi(g)$ yields 
    \begin{align}
        \pi(g)^{-1} T_G \pi(g) = \frac{1}{|G|} \sum_{h \in G} \pi(g)^{-1} \pi(h)^{-1} T \pi(h) \pi(g) = \frac{1}{|G|} \sum_{h \in G} \pi(hg)^{-1} T \pi(hg).
    \end{align}
    Since the right coset $Gg$ is equal to $G$, we can reindex the sum to obtain 
    \begin{align}
        \pi(g)^{-1} T_G \pi(g) = \frac{1}{|G|} \sum_{h \in G} \pi(hg)^{-1} T \pi(hg) = \frac{1}{|G|} \sum_{h \in G} \pi(h)^{-1} T \pi(h) = T_G.
    \end{align}
    We conclude that $\pi(g)^{-1} T_G \pi(g) = T_G$ and hence $T_G \pi(g) = \pi(g) T_G$.
\end{proof}

\subsection{Character Theory}

\Cref{item:artin-wedderburn} and \Cref{lem:twirling} show us that in order to fully understand a twirled channel, it is necessary to understand how many copies of each irreducible representation are present in a given representation.
Luckily, if the irreps are understood, then there is an efficient algorithm for computing how many copies of that irrep occur in the decomposition.
This technique comes from a subject called character theory.

\subsubsection{Characters and Conjugacy Classes}

If $\pi : G \to \mathrm{End}(V)$ is a representation, then the character $\chi_\pi : G \to \mathbb{C}$ of $\pi$ is defined by 
\begin{align}
    \chi_\pi(g) = \Tr \pi(g).
\end{align}
By the cyclic property of the trace, characters satisfy the property $\chi_\pi(h g h^{-1}) = \chi_\pi(g)$.
Because of this, characters are constant on the conjugacy classes.
Such a function is called a class function.
Evidently, the space of class functions has dimension equal to the number of conjugacy classes of $G$.
Moreover, the space is equipped with a natural inner product 
\begin{align}
    \langle \chi, \tau \rangle = \frac{1}{|G|} \sum_{g \in G} \chi(g)^* \tau(g).
\end{align}
The following lemma is among the most important results in character theory.

\begin{lemma}[Character Orthogonality \cite{sagan2001symmetric} Theorem 1.9.3] \label{lem:character-orthogonality}
    If $\pi_1$ and $\pi_2$ are irreducible representations of the finite group $G$, then 
    \begin{align}
        \langle \chi_{\pi_1}, \chi_{\pi_2}\rangle = \begin{cases}
            1 & \pi_1 \cong \pi_2, \\
            0 & \text{otherwise}.
        \end{cases}
    \end{align}
\end{lemma}

\begin{cor}[\cite{sagan2001symmetric} Corollary 1.9.4] \label{cor:character-relations}
    Let $\pi : G \to \mathrm{End}(V)$ be a representation and $\pi = m_0 \pi_1 \oplus \dots \oplus m_k \pi_k$ be the decomposition of the representation $\pi$ into irreps. The following hold.
    \begin{enumerate}
        \item The character $\chi_{\pi}$ can be written as $\chi_{\pi} = m_0 \chi_{\pi_0} + m_1 \chi_{\pi_1} + \dots + m_k \chi_{\pi_k}$.
        \item \label{item:irrep-mult-test} The inner product $\langle \chi_{\pi}, \chi_{\pi_i}\rangle = m_i$ for all $0 \leq i \leq k$.
        \item The inner product $\langle \chi_{\pi}, \chi_{\pi} \rangle = m_0^2 + \dots + m_k^2$. In particular this is always an integer.
        \item The representation $\pi$ is irreducible if and only if $\langle \chi_{\pi}, \chi_{\pi} \rangle = 1$.
        \item If $\tau$ is another representation of $G$, then $\pi$ and $\tau$ are isomorphic if and only if $\chi_{\pi}(g) = \chi_{\tau}(g)$ for all $g \in G$.
    \end{enumerate}
\end{cor}

\Cref{lem:character-orthogonality} and \Cref{cor:character-relations} show us two remarkable facts.
Firstly, the irrep characters are orthonormal and thus linearly independent.
From this we can conclude that the irrep characters form a basis for the space of class functions and, in particular, that there must be exactly as many irreps as there are conjugacy classes of $G$.
Secondly, \Cref{item:irrep-mult-test} gives us an efficient method of determining how many copies of a given irrep are present in an arbitrary representation of $G$.